\chapter{群体操作系统：AgentOS 的尺度推广与三相记忆的群体化}

\begin{chapterlead}
本章把前一章开启的尺度推广进一步推向操作系统层面。前章已经在 SFPU 的框架下讨论了流体群体智能与跨尺度具身协同，AHC 则承担了群体协调与异构协同问题；但``多个 Agent 如何通信与协作''并不是终点。当多个长期自治 Agent 形成稳定群体之后，系统面对的将是群体自身的状态、历史、势能、结构与边界——群体开始成为独立于任何单个 Agent 的认知实体，需要一个面向群体级认知状态、记忆、场结构与演化过程的操作系统层，即 SwarmOS。本章的任务，就是为这一实体建立完整的原理定义。

本章的中心论断是：AgentOS 向 SwarmOS 的尺度提升不是 $N$ 个单体系统的简单叠加，而是``相同动力学形式、不同尺度序参量''的再实例化。我们将从单个 Agent 的状态 $\mathcal{A}_i=(h_i,E_i,\Theta_i,\Psi_i,B_i)$ 出发，说明为什么 $\mathcal{S}\neq\sum_{i}\mathcal{A}_i$；继而把智能动力学统一方程与三相记忆整体平移到群体尺度——工作记忆对应群体实时状态场，情景记忆对应群体任务势能，结构记忆对应公共结构流形——并分析三相时间尺度、多尺度滤波与群体稳定性的关系。在此之上，本章沿自我与认知控制的构件逐一推广：群体自我 $\Psi_G$、群体自感知（SPC）、群体场调控（TCE）、群体复杂度（CCM）、群体认知卸载、群体边界与记忆主权，并把关系拓扑确立为 SwarmOS 的一等状态，讨论涌现序参量、群体相变，以及个体自治与群体一致性这一基本矛盾在群体记忆中的双向凝固、蒸发与重构。最后，本章给出 SwarmOS 的核心组成、SFPU 在其中的动力学计算层位置、与计算资源的层级关系及其原理定义，并把 AgentOS、SwarmOS 与 WorldOS 收束为同一多尺度智能操作系统结构。

至此，三相记忆不再只是一种记忆设计，而表现为长期自适应智能系统的普适多时间尺度动力学。完成向群体尺度的推广之后，我们把视线收回主体内部：下一章讨论三信号通道，即身体状态如何经分层路由进入认知内核。
\end{chapterlead}

前文已经分别讨论了 SFPU 所对应的群体场计算机制，以及 AHC 所承担的群体协调与异构协同问题。这里引出一个更深层的结构，当多个长期自治 Agent 形成稳定群体后，系统所面对的问题已经不再只是``多个 Agent 如何通信与协作''，而是群体本身开始形成独立于任何单个 Agent 的状态、历史、势能、结构、边界与演化规律。

因此，在 AgentOS 之上有必要引入一个新的操作系统层：

\begin{statementbox}
\centering
\adjustbox{max width=.96\linewidth}{\(\displaystyle
\mathrm{SwarmOS}
\)}
\end{statementbox}

SwarmOS 并不是简单的多智能体编排器，也不是 $N$ 个 AgentOS 的并行运行环境，而是面向群体级认知状态、群体级记忆、群体级场结构和群体级演化过程的操作系统。

AgentOS 与 SwarmOS 的关系可以概括为：

\begin{statementbox}
\centering
\adjustbox{max width=.96\linewidth}{\(\displaystyle
\mathrm{AgentOS}
\;\xrightarrow{\text{尺度提升}}\;
\mathrm{SwarmOS}
\)}
\end{statementbox}

但这种尺度提升并不是简单复制。其本质是：

\begin{statementbox}
\centering
\adjustbox{max width=.96\linewidth}{\(\displaystyle
\text{相同动力学形式}
+
\text{不同尺度序参量}
\)}
\end{statementbox}

三相记忆动力学因此不仅适用于单个 Agent，也可以被推广到群体尺度。工作记忆、情景记忆与结构记忆在群体层分别表现为群体实时状态场、群体任务势能与公共结构流形。

这意味着，三相记忆理论开始从一种 Agent 记忆模型，进一步上升为一种跨尺度智能动力学结构。


\section{从个体智能到群体智能}

单个 AgentOS 处理的是一个长期存在的智能主体。

其基本状态可以写为：

\[
\mathcal{A}_i
=
\left(
h_i,
E_i,
\Theta_i,
\Psi_i,
B_i
\right),
\]

其中：

\[
h_i(t)
\]

表示第 $i$ 个 Agent 的工作记忆状态，

\[
E_i
\]

表示其情景记忆，

\[
\Theta_i
\]

表示其结构记忆，

\[
\Psi_i
\]

表示其自我结构，

\[
B_i
\]

表示其认知边界。

对于单体系统，基本三相循环为：

\[
h_i
\leftrightarrow
E_i
\leftrightarrow
\Theta_i.
\]

当存在 $N$ 个 Agent 时，一个最直接的想法是：

\[
\mathcal{S}
=
\left\{
\mathcal{A}_1,
\mathcal{A}_2,
\ldots,
\mathcal{A}_N
\right\}.
\]

但真正的群体智能并不等于多个 Agent 的简单集合。

即：

\[
\boxed{
\mathcal{S}
\neq
\sum_{i=1}^{N}\mathcal{A}_i
}
\]

原因在于多个 Agent 发生长期耦合以后，会出现新的宏观序参量。

例如：

\[
\text{群体目标},
\]

\[
\text{群体方向},
\]

\[
\text{公共任务势能},
\]

\[
\text{公共语义空间},
\]

\[
\text{组织结构},
\]

\[
\text{角色分布},
\]

\[
\text{群体边界},
\]

\[
\text{群体身份}.
\]

这些变量不能简单归属于任何一个 Agent。

因此必须定义群体层状态：

\[
X_G(t).
\]

群体智能由此表现为：

\[
\left\{
\mathcal{A}_1,
\dots,
\mathcal{A}_N
\right\}
\xrightarrow{\mathcal{C}}
X_G,
\]

其中 $\mathcal{C}$ 表示从微观多 Agent 状态到宏观群体状态的粗粒化过程。

因此：

\[
X_G
=
\mathcal{C}
\left(
\mathcal{A}_1,
\dots,
\mathcal{A}_N,
R,
W
\right),
\]

其中 $R$ 表示 Agent 之间的关系拓扑，$W$ 表示外部世界状态。

SwarmOS 所操作的基本对象正是这个宏观群体状态，而不再只是单个 Agent 实例。


\section{智能动力学的统一方程}

个体 Agent 与群体系统可以使用同一类智能动力学形式描述：

\[
\boxed{
\dot X
=
-\nabla_g V
+
F_{\mathrm{external}}
+
\xi
}
\]

其中：

\[
X
\]

表示系统当前快速状态，

\[
V
\]

表示由历史形成的中时间尺度势能，

\[
g
\]

表示长期形成的状态空间几何，

\[
F_{\mathrm{external}}
\]

表示来自环境、身体或其他系统的外部作用，

\[
\xi
\]

表示随机扰动、探索项以及无法显式建模的残余变化。

这一方程揭示了三相记忆动力学更加一般的数学结构：

\begin{statementbox}
\centering
\adjustbox{max width=.96\linewidth}{\(\displaystyle
\text{State}
\rightarrow
\text{Potential}
\rightarrow
\text{Geometry}
\)}
\end{statementbox}

即：

当前状态首先形成历史相关势能，历史势能经过长期巩固以后进一步改变系统未来运动所处的几何空间。

因此三相记忆：

\[
h
\rightarrow
E
\rightarrow
\Theta
\]

可以更抽象地写成：

\[
X
\rightarrow
V
\rightarrow
g.
\]

对于个体 Agent：

\[
X_A = h,
\]

\[
V_A \sim E,
\]

\[
g_A \sim \Theta.
\]

于是个体智能动力学近似写为：

\[
\dot h
=
-\nabla_{\Theta}E
+
F_{\mathrm{World}}
+
F_{\Psi}
+
F_{\mathrm{Body}}
+
\xi.
\]

而对于群体系统：

\[
X_G = H_G,
\]

\[
V_G \sim E_G,
\]

\[
g_G \sim \Theta_G.
\]

于是群体智能动力学为：

\[
\boxed{
\partial_t H_G
=
-\nabla_{g_G}V_G
+
F_{\mathrm{environment}}
+
F_{\mathrm{agents}}
+
F_{\Psi_G}
+
F_R
+
\xi_G
}
\]

其中 $F_R$ 表示群体关系拓扑带来的作用。

由此可见：

\begin{statementbox}
\centering
\adjustbox{max width=.96\linewidth}{\(\displaystyle
\text{AgentOS 与 SwarmOS 共享同一类动力学母方程}
\)}
\end{statementbox}

区别不在于基本形式，而在于状态变量所代表的尺度。


\section{三相记忆的群体推广}

个体三相记忆为：

\[
h
\leftrightarrow
E
\leftrightarrow
\Theta.
\]

在群体尺度上，可以相应定义：

\[
\boxed{
H_G
\leftrightarrow
E_G
\leftrightarrow
\Theta_G
}
\]

其中：

\[
H_G
\]

为群体工作状态场，

\[
E_G
\]

为群体情景记忆或任务势能，

\[
\Theta_G
\]

为群体结构记忆或公共流形。

这三者不是三个数据库，而是三个不同时间尺度上的群体状态。

群体工作记忆 $H_G$ 对应群体当前正在发生的状态。

例如：

\[
\text{当前目标分布},
\]

\[
\text{个体位置},
\]

\[
\text{当前资源分配},
\]

\[
\text{局部冲突},
\]

\[
\text{通信状态},
\]

\[
\text{即时任务方向}.
\]

其变化速度最快。

群体情景记忆 $E_G$ 表示一定时间范围内由任务历史形成的动态结构。

例如：

\[
\text{反复成功的合作方式},
\]

\[
\text{某一区域的风险积累},
\]

\[
\text{任务阶段形成的局部吸引子},
\]

\[
\text{某种资源分配造成的历史偏置}.
\]

它不是瞬时状态，也不是永久结构，而是一种具有历史惯性的中间态。

群体结构记忆 $\Theta_G$ 则描述长期稳定的群体生成结构。

例如：

\[
\text{公共语义空间},
\]

\[
\text{长期组织规则},
\]

\[
\text{共同任务坐标},
\]

\[
\text{协同几何},
\]

\[
\text{长期角色分工},
\]

\[
\text{公共能力结构}.
\]

因此：

\[
\boxed{
H_G
\rightarrow
E_G
\rightarrow
\Theta_G
}
\]

表示群体经历从即时状态逐步转化为公共经验，再逐步转化为群体长期结构的过程。


\section{SFPU 中的三相记忆结构}

在 SFPU 的群体场处理中，已经可以观察到这种三相结构。

快速变化的是目标嵌入场：

\[
Z^\star(x,t),
\]

以及：

\[
T(t),
\qquad
\eta(t),
\qquad
\kappa(t).
\]

这些变量描述群体当前正在朝向什么方向、如何响应局部扰动，以及当前的协调状态。

因此可以定义：

\[
\boxed{
H_G
\sim
Z^\star(x,t)
}
\]

即目标场构成群体工作记忆的一种场表达。

中速更新的是任务势能与适配结构：

\[
V_\phi,
\qquad
F_\psi,
\qquad
A_{\mathrm{task}}.
\]

这些变量由一段时间中的任务经验逐渐形成，并持续影响之后的群体运动。

因此：

\[
\boxed{
E_G
\sim
\left\{
V_\phi,
F_\psi,
A_{\mathrm{task}}
\right\}
}
\]

它们构成群体情景记忆的动力学表达。

而更新最慢的是公共流形相关变量：

\[
E_S,
\qquad
E_G,
\qquad
g_Z.
\]

这些变量改变的不只是当前任务，而是群体共同使用的坐标、关系和状态空间几何。

因此：

\[
\boxed{
\Theta_G
\sim
\left\{
E_S,
E_G,
g_Z
\right\}
}
\]

其中尤其值得注意的是 $g_Z$。

它所表达的并不是一条记忆内容，而是群体状态空间本身的几何。

如果：

\[
g_Z
\rightarrow
g_Z',
\]

那么即使面对相同目标，群体最容易产生的路径也会发生改变。

因此，群体结构学习的本质不是增加更多存储内容，而是：

\begin{statementbox}
\centering
\adjustbox{max width=.96\linewidth}{\(\displaystyle
\text{改变未来群体行为的生成几何}
\)}
\end{statementbox}

这与 AgentOS 中 $\Theta$ 的定义完全一致。


\section{三相时间尺度与群体稳定性}

单体三相记忆具有显著不同的更新速度：

\[
|\dot{\Theta}|
\ll
|\dot E|
\ll
|\dot h|.
\]

群体系统同样要求：

\[
\boxed{
|\dot{\Theta}_G|
\ll
|\dot E_G|
\ll
|\dot H_G|
}
\]

SFPU 中对应为：

\[
|\dot{\theta}_{\mathrm{backbone}}|
\ll
|\dot{\phi}_{\mathrm{energy}}|
\ll
|\partial_t Z^\star|.
\]

这并不是单纯为了工程方便，而是群体长期稳定的基本条件之一。

定义三个特征时间尺度：

\[
\tau_H,
\qquad
\tau_E,
\qquad
\tau_\Theta,
\]

则必须满足：

\[
\boxed{
\tau_H
\ll
\tau_E
\ll
\tau_\Theta
}
\]

如果：

\[
\tau_\Theta
\approx
\tau_H,
\]

则任何瞬时扰动都会快速改变公共结构。

系统将出现：

\[
\text{结构漂移},
\]

\[
\text{公共坐标不稳定},
\]

\[
\text{群体行为不可预测}.
\]

反之，如果：

\[
\tau_\Theta
\rightarrow
\infty,
\]

则群体永久无法改变其公共结构，又会失去适应性。

因此三相结构解决的是长期智能系统的一个根本矛盾：

\begin{statementbox}
\centering
\adjustbox{max width=.96\linewidth}{\(\displaystyle
\text{快速适应}
\quad\text{与}\quad
\text{长期稳定}
\)}
\end{statementbox}

之间的统一。


\section{三相记忆作为多尺度滤波过程}

三相系统也可以理解为一种多时间尺度滤波结构。

设外部扰动为：

\[
u(t).
\]

快速工作态满足：

\[
\tau_H \dot H_G
=
F_H(u,H_G,E_G,\Theta_G).
\]

群体情景态相当于对工作态进行中尺度积累：

\[
\tau_E \dot E_G
=
F_E(H_G,E_G,\Theta_G).
\]

长期结构则进一步对群体经验进行低频提取：

\[
\tau_\Theta \dot \Theta_G
=
F_\Theta(E_G,\Theta_G).
\]

满足：

\[
\tau_H
\ll
\tau_E
\ll
\tau_\Theta.
\]

定义：

\[
\epsilon_1
=
\frac{\tau_H}{\tau_E},
\]

\[
\epsilon_2
=
\frac{\tau_E}{\tau_\Theta},
\]

则：

\[
\epsilon_1
\ll
1,
\]

\[
\epsilon_2
\ll
1.
\]

于是系统可以近似写成：

\[
\dot H_G
=
F_H(H_G;E_G,\Theta_G),
\]

\[
\dot E_G
=
\epsilon_1
F_E(\langle H_G\rangle,E_G,\Theta_G),
\]

\[
\dot\Theta_G
=
\epsilon_1\epsilon_2
F_\Theta(\langle E_G\rangle,\Theta_G).
\]

这里：

\[
\langle H_G\rangle
\]

表示对快速群体状态进行时间粗粒化后的统计结构。

因此群体学习本质上是：

\begin{statementbox}
\centering
\adjustbox{max width=.96\linewidth}{\(\displaystyle
\text{即时状态}
\rightarrow
\text{历史统计}
\rightarrow
\text{长期几何}
\)}
\end{statementbox}


\section{SwarmOS 不是多个 AgentOS 的简单叠加}

虽然 SwarmOS 建立在 AgentOS 之上，但：

\begin{statementbox}
\centering
\adjustbox{max width=.96\linewidth}{\(\displaystyle
\mathrm{SwarmOS}
\neq
N\times\mathrm{AgentOS}
\)}
\end{statementbox}

原因在于群体层存在不可约的宏观变量。

例如定义群体状态：

\[
H_G
=
\mathcal{C}
\left(
h_1,
h_2,
\dots,
h_N,
R,
W
\right),
\]

一般有：

\[
H_G
\neq
\sum_{i=1}^{N}h_i.
\]

这是因为群体状态不仅由个体状态决定，也取决于：

\[
R,
\]

即个体之间如何连接；

以及：

\[
W,
\]

即这些个体共同嵌入的环境。

同样：

\[
E_G
\neq
\sum_i E_i,
\]

因为群体任务历史包含大量关系性经验。

例如：

\[
A_i
\text{ 与 }
A_j
\text{ 如何协作}
\]

并不能归入 $E_i$ 或 $E_j$ 中任意一方。

更进一步：

\[
\Theta_G
\neq
\sum_i\Theta_i.
\]

因为组织规则、公共语义坐标和共同结构往往只存在于群体层。

因此群体三相记忆具有独立存在的必要性。


\section{从微观 Agent 状态到宏观群体场}

个体到群体的转换可以看成一个粗粒化过程。

设单体状态：

\[
X_i(t)
\]

满足：

\[
\dot X_i
=
-\nabla_{g_i}V_i
+
F_i
+
\sum_j C_{ij}
+
\xi_i,
\]

其中：

\[
C_{ij}
\]

表示 Agent $j$ 对 Agent $i$ 的耦合作用。

群体状态由：

\[
X_G
=
\mathcal{C}
\left(
X_1,
\dots,
X_N,
R
\right)
\]

得到。

当大量局部交互形成稳定统计结构以后，可以得到宏观场方程：

\[
\partial_t X_G
=
-\nabla_{g_G}V_G
+
F_G
+
\xi_G.
\]

这一过程类似于：

\begin{statementbox}
\centering
\adjustbox{max width=.96\linewidth}{\(\displaystyle
\text{Micro State}
\rightarrow
\text{Coarse Graining}
\rightarrow
\text{Macro Field}
\)}
\end{statementbox}

但群体系统并不是单向的。

宏观群体场形成以后，又会反过来影响个体：

\[
X_G
\rightarrow
F_i.
\]

因此完整结构为：

\begin{statementbox}
\centering
\adjustbox{max width=.96\linewidth}{\(\displaystyle
\text{Individual}
\rightarrow
\text{Swarm}
\rightarrow
\text{Individual}
\)}
\end{statementbox}

这构成 AgentOS 与 SwarmOS 之间最基本的跨尺度闭环。


\section{群体自我 $\Psi_G$}

AgentOS 中存在：

\[
\Psi_i,
\]

它使得个体当前状态并不是纯粹由外界信息推动，而始终受到``我是谁''这一自我结构的作用。

群体层同样可以出现：

\[
\boxed{
\Psi_G
}
\]

即群体自我。

它回答的问题不再是：

\[
\text{我是谁？}
\]

而是：

\begin{statementbox}
\centering
\adjustbox{max width=.96\linewidth}{\(\displaystyle
\text{我们是谁？}
\)}
\end{statementbox}

例如一个长期工作的机器人建设群体可能逐渐形成：

\[
\text{共同目标},
\]

\[
\text{共同身份},
\]

\[
\text{公共价值优先级},
\]

\[
\text{长期组织目标},
\]

\[
\text{共同风险偏好}.
\]

这些结构共同形成：

\[
\Psi_G.
\]

于是群体动力学应进一步写为：

\[
\partial_t H_G
=
-\nabla_{g_G}V_G
+
F_{\Psi_G}
+
F_{\mathrm{environment}}
+
F_R
+
\xi_G.
\]

群体自我并不要求群体具有与人类完全相同的主观意识。

在工程上，它首先是一种：

\begin{statementbox}
\centering
\adjustbox{max width=.96\linewidth}{\(\displaystyle
\text{稳定的群体自参考结构}
\)}
\end{statementbox}

即群体能够使用自身整体状态作为下一步决策的重要输入。


\section{SPC 从个体自感知到群体自感知}

AgentOS 中 SPC 的基本功能是：

\[
SPC:
h
\rightarrow
\hat h_{\Psi}.
\]

即按照自我方向对当前工作记忆状态进行压缩。

其结果重新进入：

\[
h,
\]

从而形成：

\[
h
\rightarrow
SPC
\rightarrow
\hat h_{\Psi}
\rightarrow
h'.
\]

SwarmOS 可以定义相应的群体 SPC：

\[
\boxed{
SPC_G:
H_G
\rightarrow
\hat H_{\Psi_G}
}
\]

其作用是：

\[
\text{将复杂群体场压缩为群体自身可使用的整体状态描述}.
\]

例如群体 SPC 可以输出：

\[
\text{目标一致度},
\]

\[
\text{群体分裂度},
\]

\[
\text{局部拥堵程度},
\]

\[
\text{资源压力},
\]

\[
\text{内部冲突水平},
\]

\[
\text{任务进展状态},
\]

\[
\text{组织稳定性}.
\]

这样，群体便获得一种：

\begin{statementbox}
\centering
\adjustbox{max width=.96\linewidth}{\(\displaystyle
\text{我们当前处于什么状态}
\)}
\end{statementbox}

的实时表征。

因此 Group SPC 是 SwarmOS 从普通多智能体系统走向群体认知系统的关键结构之一。


\section{TCE 从个体认知调控到群体场调控}

AgentOS 中 TCE 调节的是：

\[
T,
\]

以及增益、抑制等全局认知参数。

其本质是改变工作记忆中信息传播与状态迁移的方式。

在 SwarmOS 中，对应定义：

\[
T_G,
\]

\[
\eta_G,
\]

\[
\kappa_G,
\]

以及一个尤其重要的新参数：

\[
\lambda_G,
\]

其中 $\lambda_G$ 表示群体耦合强度。

当：

\[
\lambda_G
\uparrow,
\]

个体行为更加受到公共场影响。

群体表现为：

\[
\text{高一致性},
\]

\[
\text{高协同},
\]

\[
\text{低个体自由度}.
\]

当：

\[
\lambda_G
\downarrow,
\]

个体获得更多局部自主性。

系统表现为：

\[
\text{更高探索性},
\]

\[
\text{更强异质性},
\]

\[
\text{更弱全局一致性}.
\]

因此 SwarmOS 的 TCE 不只是调节推理温度，而是在调节：

\begin{statementbox}
\centering
\adjustbox{max width=.96\linewidth}{\(\displaystyle
\text{个体自治}
\leftrightarrow
\text{群体一致}
\)}
\end{statementbox}

之间的动态平衡。


\section{CCM 从工作记忆复杂度推广到群体复杂度}

AgentOS 中 CCM 的基本约束是：

\[
C(h)
\leq
C_{\max}.
\]

即工作记忆中的认知复杂度不能无限增长。

群体系统同样存在容量限制。

但群体复杂度不再只是 token 数量或信息量，而至少与以下因素相关：

\[
N,
\]

Agent 数量；

\[
|R|,
\]

关系数量；

\[
K,
\]

并行任务数；

\[
D,
\]

任务依赖深度；

\[
U,
\]

不确定性；

\[
\Gamma,
\]

冲突程度；

\[
L_c,
\]

通信负载。

因此可以定义：

\[
C_G
=
f
\left(
N,
|R|,
K,
D,
U,
\Gamma,
L_c
\right).
\]

SwarmOS 必须保持：

\[
\boxed{
C_G
\leq
C_{G,\max}
}
\]

当系统接近：

\[
C_{G,\max},
\]

群体 CCM 可以采取：

\[
\text{分群},
\]

\[
\text{层级化},
\]

\[
\text{创建子 Swarm},
\]

\[
\text{降低通信频率},
\]

\[
\text{压缩公共状态},
\]

\[
\text{减少全局同步},
\]

\[
\text{提高局部自治},
\]

\[
\text{将任务卸载至环境或工具}.
\]

因此群体 CCM 是组织结构形成的一个直接动力来源。

从这一角度看，层级组织并不只是人为设计结果，也可以被理解为：

\begin{statementbox}
\centering
\adjustbox{max width=.96\linewidth}{\(\displaystyle
\text{群体复杂度超过平面协调能力后的自然卸载结构}
\)}
\end{statementbox}


\section{群体认知卸载}

AgentOS 中认知卸载表示：

\[
\text{Internal Cognition}
\rightarrow
M_e.
\]

即将原本需要内部持续维持的信息、程序和能力迁移到外部结构。

SwarmOS 中，卸载出现三个尺度。

第一类是 Agent 间卸载。

例如 Agent $i$ 不再保存某项完整能力，而依赖 Agent $j$：

\[
\Theta_i
\rightarrow
Agent_j.
\]

此时 Agent $i$ 只需要保留：

\[
\text{如何找到并调用 }Agent_j.
\]

这种机制可以定义为：

\begin{statementbox}
\centering
\adjustbox{max width=.96\linewidth}{\(\displaystyle
\text{Social Offloading}
\)}
\end{statementbox}

第二类是群体共享外部记忆：

\[
E_G,\Theta_G
\rightarrow
M_{e,G}.
\]

例如：

\[
\text{共享地图},
\]

\[
\text{共享知识库},
\]

\[
\text{公共任务系统},
\]

\[
\text{标准流程}.
\]

第三类是环境卸载。

群体可以直接修改环境，使环境本身承载认知结构。

例如：

\[
\text{道路},
\]

\[
\text{仓库布局},
\]

\[
\text{标识系统},
\]

\[
\text{建筑结构},
\]

\[
\text{工具部署}.
\]

于是：

\begin{statementbox}
\centering
\adjustbox{max width=.96\linewidth}{\(\displaystyle
\text{Environment}
=
\text{Externalized Group Memory}
\)}
\end{statementbox}

群体越成熟，其智能越不应仅仅存在于单个 Agent 内部，而会分布于：

\[
\text{Agents}
+
\text{Relations}
+
\text{Tools}
+
\text{Infrastructure}
+
\text{Environment}.
\]


\section{群体边界与记忆主权}

AgentOS 中：

\[
B_i
\]

决定什么属于智能体内部，什么属于外部。

SwarmOS 则至少具有两层边界：

\[
B_i
\]

表示个体边界，

\[
B_G
\]

表示群体边界。

因此系统结构为：

\[
\text{World}
\rightarrow
B_G
\rightarrow
\text{Swarm}
\rightarrow
B_i
\rightarrow
Agent_i.
\]

这产生一个 AgentOS 中并不突出的新问题：

\begin{statementbox}
\centering
\adjustbox{max width=.96\linewidth}{\(\displaystyle
\text{什么信息属于个体，什么信息属于群体？}
\)}
\end{statementbox}

并进一步形成：

\begin{statementbox}
\centering
\adjustbox{max width=.96\linewidth}{\(\displaystyle
\text{Memory Sovereignty}
\)}
\end{statementbox}

即记忆主权问题。

例如，某段经历首先存在于：

\[
E_i.
\]

只有当它具有足够高的群体价值时，才应该：

\[
E_i
\rightarrow
E_G.
\]

而某些群体结构又可能需要被下沉至个体：

\[
\Theta_G
\rightarrow
\Theta_i.
\]

因此群体记忆流转不仅包含：

\[
H_G
\rightarrow
E_G
\rightarrow
\Theta_G,
\]

还包括：

\[
E_i
\leftrightarrow
E_G,
\]

\[
\Theta_i
\leftrightarrow
\Theta_G.
\]

SwarmOS 必须决定这些跨边界迁移何时发生。


\section{关系拓扑作为 SwarmOS 的一等状态}

AgentOS 中，一个个体的主要变化发生在：

\[
h,
\quad
E,
\quad
\Theta.
\]

而在 SwarmOS 中，除了：

\[
H_G,
\quad
E_G,
\quad
\Theta_G
\]

以外，还必须单独考虑：

\[
\boxed{
R_G(t)
}
\]

即群体关系拓扑。

这是因为群体系统不仅会改变状态，还会改变``谁和谁发生作用''。

例如：

\[
R_G(t)
\rightarrow
R_G(t+\Delta t)
\]

可能代表：

\[
\text{成员重组},
\]

\[
\text{角色变化},
\]

\[
\text{子群形成},
\]

\[
\text{指挥链变化},
\]

\[
\text{通信路径调整}.
\]

因此 SwarmOS 更完整的状态应该写成：

\[
\boxed{
\mathcal{S}_G
=
\left(
H_G,
E_G,
\Theta_G,
R_G,
\Psi_G,
B_G
\right)
}
\]

而拓扑本身满足：

\[
\dot R_G
=
F_R
\left(
H_G,
E_G,
\Theta_G,
R_G
\right).
\]

同时 $R_G$ 又反过来进入群体状态传播过程：

\[
\partial_tH_G
=
F_H
\left(
H_G,
E_G,
\Theta_G,
R_G
\right).
\]

因此：

\begin{statementbox}
\centering
\adjustbox{max width=.96\linewidth}{\(\displaystyle
\text{在群体智能中，连接方式本身就是计算的一部分}
\)}
\end{statementbox}


\section{涌现序参量}

SwarmOS 还需要管理 AgentOS 中不存在或不显著的宏观序参量。

可以定义：

\[
m_{\mathrm{align}}
\]

表示群体方向一致程度，

\[
m_{\mathrm{cohesion}}
\]

表示群体凝聚程度，

\[
m_{\mathrm{entropy}}
\]

表示群体行为熵，

\[
m_{\mathrm{fragment}}
\]

表示群体分裂程度，

\[
m_{\mathrm{trust}}
\]

表示群体信任结构，

\[
m_{\mathrm{load}}
\]

表示整体协调负载。

于是：

\[
H_G
=
\left(
Z_G,
m_{\mathrm{align}},
m_{\mathrm{cohesion}},
m_{\mathrm{entropy}},
m_{\mathrm{fragment}},
\dots
\right).
\]

这些序参量并不存在于单个 Agent 中。

它们是：

\begin{statementbox}
\centering
\adjustbox{max width=.96\linewidth}{\(\displaystyle
\text{多主体相互作用后的宏观涌现变量}
\)}
\end{statementbox}

SwarmOS 的一个重要任务就是持续估计和调控这些变量。


\section{群体认知状态与相变}

AgentOS 中工作记忆存在不同运行相态：

\[
\text{基态},
\]

\[
\text{湍流态},
\]

\[
\text{临界态},
\]

\[
\text{心流态},
\]

\[
\text{气态}.
\]

同样的思想可以推广到群体。

群体基态表示：

\[
\text{任务稳定},
\quad
\text{通信正常},
\quad
\text{组织结构稳定}.
\]

群体心流态表示：

\[
\text{目标高度一致},
\]

\[
\text{局部自治充分},
\]

\[
\text{通信负载较低},
\]

\[
\text{整体任务效率极高}.
\]

群体湍流态表现为：

\[
\text{反复冲突},
\]

\[
\text{资源争夺},
\]

\[
\text{任务互相阻塞},
\]

\[
\text{局部指令传播混乱}.
\]

群体临界态意味着：

\[
\text{组织即将发生结构转变}.
\]

例如：

\[
\text{分裂},
\]

\[
\text{合并},
\]

\[
\text{角色重组},
\]

\[
\text{任务策略突变}.
\]

群体气态则表现为：

\[
\text{高自由度},
\]

\[
\text{低一致性},
\]

\[
\text{高探索},
\]

\[
\text{弱结构约束}.
\]

因此 Group SPC 可以识别：

\[
State_G
\in
\{
\text{Base},
\text{Flow},
\text{Turbulent},
\text{Critical},
\text{Gas}
\},
\]

而 Group TCE 根据状态进行迁移控制。


\section{个体自治与群体一致性的基本矛盾}

SwarmOS 比 AgentOS 多出的一个核心问题是：

\begin{statementbox}
\centering
\adjustbox{max width=.96\linewidth}{\(\displaystyle
\text{Individual Autonomy}
\leftrightarrow
\text{Collective Coherence}
\)}
\end{statementbox}

在 AgentOS 内部，不同认知状态最终仍然服务于一个主体：

\[
\Psi_i.
\]

而群体中同时存在：

\[
\Psi_i
\]

和：

\[
\Psi_G.
\]

因此可能出现：

\[
Goal_i
\neq
Goal_G.
\]

如果群体控制过强：

\[
\lambda_G
\rightarrow
\infty,
\]

那么：

\[
Agent_i
\rightarrow
\text{纯执行单元}.
\]

群体虽然高度一致，却失去了：

\[
\text{局部适应},
\]

\[
\text{并行探索},
\]

\[
\text{多样性}.
\]

反之，如果：

\[
\lambda_G
\rightarrow
0,
\]

则系统退化成：

\[
N
\]

个相互弱关联的独立 Agent。

因此 SwarmOS 的一个核心控制目标可以写成：

\[
\max
\left[
U_G
+
\alpha
\sum_iU_i
\right]
\]

并满足：

\[
C_G
\leq
C_{G,\max}.
\]

其中：

\[
U_G
\]

表示群体整体效用，

\[
U_i
\]

表示个体效用，

\[
\alpha
\]

表示对个体自治性的保留权重。

这实际上构成了 SwarmOS 的一个基本调度问题。


\section{群体三相记忆中的双向凝固}

单体记忆中的主要凝固方向为：

\[
h
\rightarrow
E
\rightarrow
\Theta.
\]

群体系统中则存在两种不同凝固过程。

第一种是群体内部纵向凝固：

\[
H_G
\rightarrow
E_G
\rightarrow
\Theta_G.
\]

第二种是跨尺度凝固：

\[
h_i
\rightarrow
H_G,
\]

\[
E_i
\rightarrow
E_G,
\]

\[
\Theta_i
\rightarrow
\Theta_G.
\]

例如某个 Agent 发现了一条稳定规律。

初始只存在于：

\[
E_i.
\]

经过验证以后，它可能进入：

\[
E_G.
\]

进一步长期验证以后：

\[
E_G
\rightarrow
\Theta_G.
\]

此后公共结构又可能下沉：

\[
\Theta_G
\rightarrow
\Theta_j.
\]

因此群体学习形成如下循环：

\begin{statementbox}
\centering
\adjustbox{max width=.96\linewidth}{\(\displaystyle
\text{Individual Experience}
\rightarrow
\text{Group Experience}
\rightarrow
\text{Group Structure}
\rightarrow
\text{Individual Structure}
\)}
\end{statementbox}

这使得群体具备一种单体系统无法实现的学习加速机制。


\section{群体记忆中的蒸发、凝固与重构}

群体记忆同样不应该被看成永久累积的数据库。

对于群体工作态：

\[
H_G
\]

其大量状态会快速消失。

对于群体情景记忆：

\[
E_G
\]

则存在慢速蒸发：

\[
\frac{dE_G}{dt}
=
-\lambda_EE_G
+
F_E(H_G).
\]

其中：

\[
\lambda_E
\]

控制群体历史痕迹衰减。

只有重复出现、具有高价值或高预测力的结构才会：

\[
E_G
\rightarrow
\Theta_G.
\]

而 $\Theta_G$ 本身也不是绝对静态。

长期存在的新证据可能产生：

\[
\Theta_G
\rightarrow
\Theta_G'.
\]

因此群体演化仍然保持：

\begin{statementbox}
\centering
\adjustbox{max width=.96\linewidth}{\(\displaystyle
\text{快速流动}
+
\text{中速历史}
+
\text{慢速重构}
\)}
\end{statementbox}

三种过程。


\section{SwarmOS 的核心组成}

基于 AgentOS 的结构推广，SwarmOS 至少应包含以下核心机制。

群体工作状态场：

\[
H_G.
\]

群体情景势能：

\[
E_G.
\]

群体结构记忆：

\[
\Theta_G.
\]

群体自我：

\[
\Psi_G.
\]

关系拓扑：

\[
R_G.
\]

群体边界：

\[
B_G.
\]

群体自感知压缩器：

\[
SPC_G.
\]

群体温度与控制引擎：

\[
TCE_G.
\]

群体认知复杂度管理：

\[
CCM_G.
\]

群体卸载系统：

\[
Offloading_G.
\]

因此可写为：

\begin{statementbox}
\centering
\adjustbox{max width=.96\linewidth}{\(\displaystyle
\mathrm{SwarmOS}
=
\{
H_G,
E_G,
\Theta_G,
\Psi_G,
R_G,
B_G,
SPC_G,
TCE_G,
CCM_G,
Offloading_G
\}
\)}
\end{statementbox}

这构成 SwarmOS 原理层的核心状态空间。


\section{SFPU 在 SwarmOS 中的位置}

SFPU 不应被简单理解为多机器人系统中的一个普通硬件加速器。

如果 SwarmOS 的核心计算对象是：

\[
H_G,
\]

\[
E_G,
\]

\[
\Theta_G,
\]

以及：

\[
R_G,
\]

那么 SFPU 的主要职责就是加速这些群体场和群体结构的演化。

因此可以定义：

\begin{statementbox}
\centering
\adjustbox{max width=.96\linewidth}{\(\displaystyle
\mathrm{SFPU}
=
\mathrm{Swarm\ Field\ Processing\ Unit}
\)}
\end{statementbox}

其核心计算可以抽象为：

\[
(H_G,E_G,\Theta_G,R_G)
\rightarrow
(H_G',E_G',\Theta_G',R_G').
\]

传统 CPU 的原生对象主要是：

\[
\text{Instruction}.
\]

GPU 的原生对象主要是：

\[
\text{Tensor / Parallel Data}.
\]

而 SFPU 所面对的原生对象则更加接近：

\[
\text{Field},
\]

\[
\text{Potential},
\]

\[
\text{Manifold},
\]

\[
\text{Topology}.
\]

因此 SFPU 可以被理解为：

\begin{statementbox}
\centering
\adjustbox{max width=.96\linewidth}{\(\displaystyle
\text{SwarmOS 的群体认知动力学协处理器}
\)}
\end{statementbox}

而前文 AHC 与 SFPU 所解决的异构协同、群体场计算和群体运行问题，可以在 SwarmOS 这一操作系统层获得统一解释。


\section{AgentOS、SwarmOS 与计算资源的层级关系}

由此可以形成新的操作系统层级。

首先是个体认知：

\[
\mathrm{AgentOS}.
\]

其下运行：

\[
\text{Foundation Model},
\]

\[
\text{Tool},
\]

\[
\text{Body},
\]

以及个体记忆系统。

多个 AgentOS 之上形成：

\[
\mathrm{SwarmOS}.
\]

因此：

\begin{statementbox}
\centering
\adjustbox{max width=.96\linewidth}{\(\displaystyle
\mathrm{SwarmOS}
\rightarrow
\mathrm{AgentOS}
\rightarrow
\mathrm{Model/Body}
\)}
\end{statementbox}

其含义并不是 SwarmOS 取代 AgentOS，而是：

\[
\text{SwarmOS 管理群体认知},
\]

\[
\text{AgentOS 管理个体认知}.
\]

二者分别操作不同尺度的状态。


\section{AgentOS 原理向 SwarmOS 的平移}

AgentOS 中的大量核心机制可以通过尺度推广得到 SwarmOS 版本。

其对应关系为：

\[
h
\rightarrow
H_G,
\]

\[
E
\rightarrow
E_G,
\]

\[
\Theta
\rightarrow
\Theta_G,
\]

\[
\Psi
\rightarrow
\Psi_G,
\]

\[
SPC
\rightarrow
SPC_G,
\]

\[
TCE
\rightarrow
TCE_G,
\]

\[
CCM
\rightarrow
CCM_G,
\]

\[
Boundary
\rightarrow
B_G,
\]

\[
Offloading
\rightarrow
Offloading_G.
\]

但是这种推广还必须增加群体特有机制：

\[
R_G,
\]

即关系拓扑；

\[
M_G,
\]

即宏观序参量；

以及：

\[
Authority_G,
\]

即个体与群体之间的控制权分配机制。

因此可以写为：

\begin{statementbox}
\centering
\adjustbox{max width=.96\linewidth}{\(\displaystyle
\mathrm{SwarmOS}
=
\mathrm{ScaleUp}(\mathrm{AgentOS})
+
R_G
+
M_G
+
Authority_G
\)}
\end{statementbox}

这准确表达了两者既同构又不同的关系。


\section{从三相记忆到三相智能动力学}

单体 AgentOS 与群体 SwarmOS 得到统一以后，可以重新理解三相记忆。

最初我们将其描述为：

\[
h
\rightarrow
E
\rightarrow
\Theta.
\]

从更一般的动力学角度，它实际上对应：

\[
\boxed{
X
\rightarrow
V
\rightarrow
g
}
\]

其中：

\[
X
\]

是当前状态，

\[
V
\]

是历史形成的势能，

\[
g
\]

是长期形成的生成几何。

因此：

\[
\boxed{
\dot X
=
-\nabla_gV
+
F_{\mathrm{external}}
+
\xi
}
\]

可以被理解为三相智能动力学的基础运动方程。

中尺度和慢尺度则满足：

\[
\dot V
=
\epsilon
\mathcal{F}_V(X,V,g),
\]

\[
\dot g
=
\epsilon\delta
\mathcal{F}_g(V,g),
\]

并满足：

\[
1
\gg
\epsilon
\gg
\epsilon\delta.
\]

于是整个系统形成：

\begin{statementbox}
\centering
\adjustbox{max width=.96\linewidth}{\(\displaystyle
\begin{aligned}
\text{Fast:}\quad
&X
&&\text{--- 当前状态},\\
\text{Medium:}\quad
&V
&&\text{--- 历史势能},\\
\text{Slow:}\quad
&g
&&\text{--- 生成几何}.
\end{aligned}
\)}
\end{statementbox}

这比单纯将三相定义为三种记忆存储形式更加一般。

三相记忆由此可以被重新定义为：

\begin{statementbox}
\centering
\adjustbox{max width=.96\linewidth}{\(\displaystyle
\text{智能系统在多时间尺度下，
把瞬时状态逐步转化为历史势能，
再进一步转化为生成几何的动力学过程}
\)}
\end{statementbox}


\section{AgentOS、SwarmOS 与 WorldOS 的统一尺度}

群体尺度进一步向上，还可以进入世界尺度。

AgentOS 管理：

\begin{statementbox}
\centering
\adjustbox{max width=.96\linewidth}{\(\displaystyle
\text{个体认知动力学}
\)}
\end{statementbox}

SwarmOS 管理：

\begin{statementbox}
\centering
\adjustbox{max width=.96\linewidth}{\(\displaystyle
\text{群体认知动力学}
\)}
\end{statementbox}

WorldOS 管理：

\begin{statementbox}
\centering
\adjustbox{max width=.96\linewidth}{\(\displaystyle
\text{世界状态与长期世界演化动力学}
\)}
\end{statementbox}

于是形成：

\begin{statementbox}
\centering
\adjustbox{max width=.96\linewidth}{\(\displaystyle
\mathrm{AgentOS}
\rightarrow
\mathrm{SwarmOS}
\rightarrow
\mathrm{WorldOS}
\)}
\end{statementbox}

三个尺度分别具有自己的：

\[
X,
\]

\[
V,
\]

\[
g.
\]

在 AgentOS 中：

\[
X=h,
\qquad
V=E,
\qquad
g=\Theta.
\]

在 SwarmOS 中：

\[
X=H_G,
\qquad
V=E_G,
\qquad
g=\Theta_G.
\]

在 WorldOS 中，也可以相应存在：

\[
X_W,
\qquad
V_W,
\qquad
g_W,
\]

分别描述当前世界状态、历史世界势能和长期世界结构。

因此三个 OS 并不是彼此无关的软件系统，而可以被理解成：

\begin{statementbox}
\centering
\adjustbox{max width=.96\linewidth}{\(\displaystyle
\text{同一种多尺度智能动力学在不同空间与主体尺度上的操作系统实现}
\)}
\end{statementbox}


\section{SwarmOS 的原理定义}

基于上述讨论，可以对 SwarmOS 给出一个统一定义。

\begin{statementbox}
SwarmOS 是建立在群体三相记忆动力学之上的群体认知操作系统。它以群体工作状态场 $H_G$ 表示实时群体认知，以群体情景记忆 $E_G$ 表示历史形成的任务势能，以群体结构记忆 $\Theta_G$ 表示长期公共生成几何，并进一步通过群体自我 $\Psi_G$、关系拓扑 $R_G$、群体 SPC、TCE、CCM、Boundary 与 Offloading，实现个体自治、群体协调、群体记忆巩固、组织拓扑演化以及跨尺度认知调度。
\end{statementbox}

因此：

\begin{statementbox}
\centering
\adjustbox{max width=.96\linewidth}{\(\displaystyle
\mathrm{SwarmOS}
\neq
\text{Multi-Agent Framework}
\)}
\end{statementbox}

多智能体框架主要解决：

\[
\text{多个 Agent 如何通信、分工和调用}.
\]

SwarmOS 所解决的是：

\begin{statementbox}
\centering
\adjustbox{max width=.96\linewidth}{\(\displaystyle
\text{一个由多个自治智能体组成的群体，
如何形成自身的状态、历史、结构、自我和生命周期}
\)}
\end{statementbox}


\textbf{本章结论:}

AgentOS 的三相记忆并不局限于单体 Agent。

当智能主体从个体提升到群体尺度以后，同样会自然形成：

\[
\boxed{
H_G
\leftrightarrow
E_G
\leftrightarrow
\Theta_G
}
\]

其中：

\[
H_G
\]

对应群体工作记忆，

\[
E_G
\]

对应群体情景势能，

\[
\Theta_G
\]

对应群体公共结构流形。

三者仍然满足：

\[
|\dot{\Theta}_G|
\ll
|\dot E_G|
\ll
|\dot H_G|.
\]

AgentOS 与 SwarmOS 因此共享同一个更一般的智能动力学母结构：

\[
\boxed{
\dot X
=
-\nabla_gV
+
F_{\mathrm{external}}
+
\xi
}
\]

其中三相结构可以重新表达为：

\begin{statementbox}
\centering
\adjustbox{max width=.96\linewidth}{\(\displaystyle
\text{State}
\rightarrow
\text{Potential}
\rightarrow
\text{Geometry}
\)}
\end{statementbox}

AgentOS 是这一结构在个体尺度上的实现，

SwarmOS 是这一结构在群体尺度上的实现。

二者最大的区别在于：

\begin{statementbox}
\centering
\adjustbox{max width=.96\linewidth}{\(\displaystyle
\text{AgentOS 管理主体内部认知，
SwarmOS 管理主体之间形成的宏观认知场}
\)}
\end{statementbox}

SwarmOS 因此额外需要处理：

\[
\text{涌现},
\]

\[
\text{关系拓扑},
\]

\[
\text{群体自我},
\]

\[
\text{群体边界},
\]

\[
\text{控制权分配},
\]

以及：

\[
\text{个体自治与群体一致性的动态平衡}.
\]

SFPU 则可以被放置在这一理论体系的群体动力学计算层，作为：

\begin{statementbox}
\centering
\adjustbox{max width=.96\linewidth}{\(\displaystyle
\text{群体状态场、任务势能、公共流形与组织拓扑的动力学加速器}
\)}
\end{statementbox}

因此，前文关于 SFPU 与 AHC 的讨论，在 SwarmOS 中获得了一个统一的操作系统解释。

至此，AgentOS 理论从单体智能体进一步扩展到群体尺度，并形成：

\begin{statementbox}
\centering
\adjustbox{max width=.96\linewidth}{\(\displaystyle
\mathrm{AgentOS}
\rightarrow
\mathrm{SwarmOS}
\rightarrow
\mathrm{WorldOS}
\)}
\end{statementbox}

这一多尺度智能操作系统结构。

三相记忆也由此不再只是一种记忆设计，而开始表现为：

\begin{statementbox}
\centering
\adjustbox{max width=.96\linewidth}{\(\displaystyle
\text{长期自适应智能系统的普适多时间尺度动力学}
\)}
\end{statementbox}